% 2010 年安徽卷理科第 14 题：算法流程图。
% 来源 PDF：普通高考/2010/2010安徽理.pdf，SHA-256 2760eb9e09374139a51303b0c2590746c74dfef7fed99ac6b2eace305af646c5。
% 原页：第 1 页（A3 三栏）第 3 栏上部，书页码 569。
% 重绘依据：开始 → x=1 → 判断「x 是奇数?」；是（右）→ x=x+1 → 回到判断前；
%   否（下）→ x=x+2 → 判断「x>8?」；否（右）→ 回到「x 是奇数?」前；是（下）→ 输出 x → 结束。
\begin{tikzpicture}[semithick, >=Stealth,
  every node/.style={inner sep=1.5pt, font=\small}]
  \node (start) [draw, rounded corners=4pt, minimum width=1.9cm, minimum height=0.55cm] at (0,0) {开始};
  \node (x1) [draw, minimum width=1.9cm, minimum height=0.55cm] at (0,-1.35) {$x=1$};
  \node (c1) [draw, diamond, aspect=2.1, minimum width=3.1cm,
    minimum height=1.15cm] at (0,-3.15) {$x$ 是奇数?};
  \node (p1) [draw, minimum width=2.5cm, minimum height=0.55cm] at (3.7,-3.15) {$x=x+1$};
  \node (p2) [draw, minimum width=2.5cm, minimum height=0.55cm] at (0,-4.85) {$x=x+2$};
  \node (c2) [draw, diamond, aspect=2.1, minimum width=3.1cm,
    minimum height=1.15cm] at (0,-6.65) {$x>8?$};
  \node (out) [draw, trapezium, trapezium left angle=72, trapezium right angle=108,
    minimum width=2.4cm, minimum height=0.55cm] at (0,-8.25) {输出 $x$};
  \node (end) [draw, rounded corners=4pt, minimum width=1.9cm, minimum height=0.55cm] at (0,-9.45) {结束};

  \draw[->] (start) -- (x1);
  \draw[->] (x1) -- (c1);
  \draw[->] (c1.east) -- node[above] {是} (p1.west);
  \draw[->] (c1.south) -- node[right] {否} (p2.north);
  \draw[->] (p2) -- (c2);
  \draw[->] (c2.south) -- node[right] {是} (out.north);
  \draw[->] (out) -- (end);
  % x=x+1 的回流
  \draw (p1.north) -- (3.7,-2.0) -- (0,-2.0);
  % x>8? 为否时的回流
  \draw (c2.east) -- node[above] {否} (5.2,-6.65) -- (5.2,-2.0) -- (0,-2.0);
  \draw[->] (0,-2.0) -- (c1.north);
\end{tikzpicture}
